\scanpage{10}
\mspage{\unsure{8|9|1}}
d'une uniformisante \deleted{\unsure{de}} \deleted{\(\unsure{\eta|g}\)} \(t\) de \(V\) \inserted{donc de \(\widetilde V\)}, comme la réunion des sous-extensions \deleted{abéliennes} Kummeriennes
\[
K_n=\widetilde K(\sqrt[n]{t})\subset\overline K,\tag*{(2.4.4)}
\]
où \(n\) parcourt les entiers \(>0\) premiers à \(p\). Le groupe de Galois de cette extension \(K_n\) est naturellement isomorphe au groupe
\[
\mu_n(\overline K)={}_n\overline K^*
\]
des racines \(n\)ièmes de l'unité dans \(\overline K\), et par passage à la limite sur \(n\) on trouve
\[
I_t=\operatorname{Gal}((\varinjlim_n K_n)/K)\simeq\varprojlim_n\operatorname{Gal}(K_n/K),\deleted{\unsure{\mathbb Z_\ell(1)}},
\]
i.e (2.4.2).

\uline{Remarque} 2.5. On notera que par définition, le sous-groupe \(P\) de \(I\) est stable par \uline{\unsure{tout}} autom. de \(I\), en particulier par les automorphismes induits par les autom. int. de \(\pi\). Donc \(\pi\) opère par passage au quotient sur \(I/p=I_t\) ; \deleted{On constate aussitôt que} \deleted{cette \unsure{opération|op. équivaut}} ds cette opération, le ss-groupe \(I\) de \(\pi\) opère évidemment trivialement, \(I_t\) étant commutatif ; on en trouve en fait une opération de \(\pi/I\simeq\pi_0\) sur \(I_t\). \deleted{\unsure{Compte t}} On vérifie trivialement que
% archive sheet number: 9
